\scanpage{67}
Il y a aussi un foncteur fibre $T_{\mathrm{cr}}$ au \deleted{\unsure{dessus de}} dessus de $K$, à valeurs dans les vectoriels \deleted{[gradués \unsure{et}] filtrés}. Donc on a un torseur
\[
\left\{\begin{gathered}
\boxed{Q\ \text{torseur pour }G_K}\\
T_{\mathrm{cr}}\simeq Q\mathbin{\wedge}^{G_K}(T_p)_K
\end{gathered}\right.
\]
\deleted{\diagram{\[
k_2:\begin{tikzcd}
Gm_K\arrow[r]&G''=\uline{G_K^Q=\operatorname{ad}(Q)}
\end{tikzcd}
\]}
\unsure{défini modulo} conjugaison par $\in U''(K)$, $U''$ \inserted{de $G''$} le \unsure{gr} unipotent associé à $k_2$, les $k_2$ \unsure{formant} un torseur sous $U''$.}

$T_{\mathrm{cr}}\otimes_K L$ a une filtration,

(Il y a un isom. de \deleted{torseurs} foncteurs filtrés

($T_{\mathrm{Hdg}}$ filtré par $i_2$)
\[
\operatorname{gr}(T_{\unsure{\mathrm{DR}|\mathrm{cr}}}\otimes_K L)\simeq T_{\mathrm{Hdg}}
\]
qui correspond à \deleted{des isom} un torseur $R$ sous $U'\subset G'=G_L^P$
\margin{\unsure{grp unipot.}}
\deleted{\diagram{\[
\unsure{\widetilde k_2}:\begin{tikzcd}
Q_L&P\arrow[l,"\sim"]
\end{tikzcd}
\]}}
défini par $j_2$

\deleted{\unsure{satisfaisant}
\[
\widetilde{k_2u}=\unsure{k_2}u_{Q_L}\circ\widetilde k_2
\]}
\deleted{\unsure{i.e. on a un}}
\diagram{\[\boxed{\begin{tikzcd}[row sep=small]
R\arrow[r,hook]&\operatorname{Isom}(\deleted{\unsure{P}}Q_L,P)\\
\text{torseur pour }U'&\unsure{G''|G'}\supset U'
\end{tikzcd}}\]}
\margin{\deleted{il faut \unsure{exprimer} que cela \unsure{correspond} à un isom de foncteurs bigradués}}

D'autre part, il y a sur le foncteur $T_{\unsure{\mathrm{cr}|\mathrm{DR}}}$ une $F$-structure
\diagram{\[
F:\begin{tikzcd}
T_{\mathrm{DR}}^\sigma\arrow[r,"\unsure{\sim}"]&T_{\mathrm{DR}}
\end{tikzcd}
\]}
\mspage{66}
